\section{Empirical Evaluation}
\label{sec:evaluation}

We evaluated our reference implementation across three dimensions: storage overhead, execution latency, and security resilience under adversarial conditions.

\begin{table}[htbp]
\caption{Container Framing Overhead by Component}
\label{tab:overhead}
\begin{center}
\begin{tabular}{llr}
\toprule
\textbf{Component} & \textbf{Format / Encoding} & \textbf{Size (Bytes)} \\
\midrule
Fixed Header & Packed C-struct & 32 \\
Public Key & Ed25519 raw bytes & 32 \\
AEAD Nonce & AES-GCM IV & 12 \\
Digital Signature & Ed25519 signature & 64 \\
Manifest Segment & Canonical CBOR (avg) & 180--240 \\
\midrule
\textbf{Total Framing Overhead} & & \textbf{$\approx$ 320--380 Bytes} \\
\bottomrule
\end{tabular}
\end{center}
\end{table}

\subsection{Performance Latency}
On an AMD Ryzen processor, packing and signing a 10\,MB payload required an average of $8.24$\,ms. Header parsing, CBOR deserialization, and signature verification executed in $0.42$\,ms, enabling real-time scanning of hundreds of files per second without measurable UI stutter.

\subsection{Adversarial Tampering Tests}
In testing over 10,000 synthetic corruptions:
\begin{itemize}
    \item Single-bit bitflips in the expiration timestamp resulted in 100\% signature rejection.
    \item Clock rollback simulations across reboots were caught by the monotonic delta checker with zero false negatives.
\end{itemize}
